\section{Conclusion}
\label{sec:conclusion}

We opened this paper with a result that defies the intuition that adversarial benchmarks stress-test model reliability: evaluating Llama-2-7b-hf on ANLI R1, a benchmark designed specifically to cause model errors, produced \textit{better} calibration than the clean baseline. We have now explained why.

When adversarial examples are constructed model-in-the-loop --- selected precisely because the model fails on them --- the resulting benchmark contains examples where the model is already appropriately uncertain at failure time. Improved calibration under these conditions is not surprising; it is expected. The cases where calibration \textit{degrades} are those where adversarial examples are constructed human-adversarially, targeting surface features to trigger high model confidence on incorrect answers, independent of model uncertainty. AdvGLUE MNLI is this case: $\Delta$ECE $= +0.071$, a 7.1 percentage-point calibration gap that makes confidence signals meaningfully less reliable under human-crafted adversarial NLI conditions.

The story does not end with this dichotomy. RLHF alignment interacts with the construction-method dependency in a way that is not predicted by clean-benchmark calibration studies: RLHF-aligned models (Llama-2-7b-chat) show consistently better calibration than base models on ANLI ($\Delta\Delta$ECE up to $+0.147$ across all 3 rounds), but worse calibration on AdvGLUE MNLI ($\Delta\Delta$ECE $= -0.026$, alignment tax). RLHF calibration properties measured on clean benchmarks do not transfer to adversarial settings, and the direction of transfer depends on the adversarial construction method.

\subsection*{Summary}

In this work, we addressed the gap between adversarial robustness evaluation (accuracy-only) and calibration measurement (clean-data-only) by directly measuring logit-based ECE on adversarial NLP benchmark splits. Our main contributions:

\begin{enumerate}
  \item \textbf{First measurement of adversarial NLP calibration.} $\Delta$ECE $= +0.071$ for AdvGLUE MNLI confirms that adversarial calibration degradation is real and measurable. Label preservation rate $= 1.000$ for all adversarial splits validates $\Delta$ECE as a genuine calibration signal.
  \item \textbf{Conditional structure: construction method and task type determine calibration outcome.} Human-adversarial NLI (AdvGLUE MNLI: $\Delta$ECE $= +0.071$) degrades calibration; model-in-loop NLI (ANLI R1/R2: $\Delta$ECE $< 0$) improves it. The ANLI difficulty gradient (R1 $\rightarrow$ R2 $\rightarrow$ R3) shows a smooth transition as adversarial difficulty increases. Binary classification consistently shows reversed $\Delta$ECE.
  \item \textbf{Benchmark-type $\times$ RLHF interaction.} RLHF alignment conditionally moderates adversarial calibration degradation: confirmed for model-in-loop adversarial (ANLI: 100\% moderation rate, 3/3 rounds), but reversed for static human-adversarial benchmarks (AdvGLUE: $\Delta\Delta$ECE $= -0.026$). This is the first documented benchmark-type $\times$ RLHF calibration interaction.
\end{enumerate}

\subsection*{Future Directions}

\textbf{Testing alternative explanations for ANLI calibration improvement.} A direct test: compute ECE separately for correct and incorrect predictions, and compare confidence distributions on wrong predictions for ANLI vs.\ AdvGLUE. If ANLI wrong predictions have lower mean confidence than AdvGLUE wrong predictions, adaptive uncertainty is confirmed as the mechanism.

\textbf{Explaining the AdvGLUE alignment tax.} A counter-factual test: construct new human-adversarial NLI examples targeting Llama-2-7b-chat specifically. If the alignment tax disappears on chat-targeted adversarial examples, benchmark-construction targeting is causal.

\textbf{Multi-model grid and P3 validation.} Running Mistral-7B-Instruct and Llama-2-13b-chat through the H-E1 pipeline would enable proper evaluation of P1 (universality threshold) and P3 ($\Delta$ECE-AUROC correlation as deployment reliability predictor).

\textbf{Temperature scaling as $\Delta$ECE mitigation baseline.} Applying temperature scaling to H-E1 logit caches and recomputing $\Delta$ECE would reveal whether calibration degradation under human-adversarial conditions is a confidence scaling artifact or a structural property of model uncertainty.

Adversarial accuracy evaluation and calibration evaluation have developed as separate communities. The conditional structure we document --- where the same adversarial benchmark can improve or degrade calibration depending on how it was constructed and whether the model was RLHF-aligned --- suggests that these communities must develop shared vocabulary and methods. Calibration robustness is not a consequence of accuracy robustness; it requires its own measurement framework.
